\documentclass[12pt,a4paper]{article}

\usepackage[utf8]{inputenc}
\usepackage[T1]{fontenc}
\usepackage[french]{babel}
\usepackage{geometry}
\usepackage{hyperref}
\usepackage{booktabs}
\usepackage{longtable}
\usepackage{array}
\usepackage{listings}
\usepackage{xcolor}
\usepackage{titlesec}
\usepackage{parskip}

\geometry{margin=2.5cm}

% Couleurs
\definecolor{codebg}{RGB}{245,245,245}
\definecolor{codeframe}{RGB}{200,200,200}
\definecolor{primary}{RGB}{79,70,229}

% Style des titres
\titleformat{\section}{\Large\bfseries\color{primary}}{}{0em}{}[\titlerule]
\titleformat{\subsection}{\large\bfseries}{}{0em}{}

% Style des blocs de code
\lstset{
  backgroundcolor=\color{codebg},
  frame=single,
  rulecolor=\color{codeframe},
  basicstyle=\ttfamily\small,
  breaklines=true,
  keepspaces=true,
  columns=flexible,
}

\hypersetup{
  colorlinks=true,
  linkcolor=primary,
  urlcolor=primary,
}

% ─────────────────────────────────────────────
\begin{document}

\begin{titlepage}
  \centering
  \vspace*{3cm}
  {\Huge\bfseries\color{primary} Rapport — Tradrly Bot\par}
  \vspace{1cm}
  {\large Assistant IA local, multi-couches, à exécution majoritairement hors-ligne\par}
  \vspace{2cm}
  {\large\today\par}
\end{titlepage}

% ─────────────────────────────────────────────
\section{Ce que fait le projet}

Tradrly Bot est un \textbf{assistant IA local} contrôlable par voix ou texte,
qui tourne entièrement sur la machine de l'utilisateur (\emph{offline-first}).
Il comprend les commandes en français et en anglais (y compris avec fautes de frappe)
et agit directement sur le PC : ouvrir des applications, régler le volume,
contrôler la luminosité, gérer des fichiers, etc.

% ─────────────────────────────────────────────
\section{Architecture en couches}

\begin{lstlisting}
UI (Electron + React)
       | WebSocket / HTTP
Orchestrateur (Node.js :4000)   <- routeur de commandes
       | HTTP
Couche IA (Python FastAPI :8000) <- NLU / NLP local / STT / TTS
       |
Services externes (Slack, Gmail, Notion...)
\end{lstlisting}

\subsection*{Description des couches}

\begin{itemize}
  \item \textbf{UI} : interface Electron + React, affiche l'avatar animé, l'historique et les réglages. Communique avec l'orchestrateur via WebSocket.
  \item \textbf{Orchestrateur} : cœur Node.js. Reçoit les commandes, applique les règles locales (regex) et délègue à la couche IA si nécessaire.
  \item \textbf{Couche IA} : serveur Python FastAPI. Gère la transcription vocale (STT), la compréhension (NLU + NLP local via spaCy/TextBlob), la génération de réponse, la synthèse vocale (TTS) et la mémoire vectorielle.
  \item \textbf{Services} : intégrations avec des applications tierces (Slack, Teams, Discord, Notion, Gmail, Git, Calendar).
\end{itemize}

% ─────────────────────────────────────────────
\section{APIs et bibliothèques utilisées}

\begin{longtable}{>{\bfseries}p{4.5cm} p{2.5cm} p{7.5cm}}
\toprule
Bibliothèque / API & Couche & À quoi ça sert \\
\midrule
\endhead

Express & Orchestrateur & Serveur HTTP REST (\texttt{/api/command}, \texttt{/api/desktop}, \texttt{/api/history}…) \\
\addlinespace
ws (WebSocket) & Orchestrateur & Canal temps-réel entre l'UI et l'orchestrateur (état de l'avatar, réponses) \\
\addlinespace
node:sqlite (intégré Node ≥22) & Orchestrateur & Persistance de l'historique des commandes en local, sans dépendance native \\
\addlinespace
dotenv & Orchestrateur & Lecture du fichier \texttt{.env} (ports, clés API) \\
\addlinespace
cors & Orchestrateur & Autorise les requêtes cross-origin de l'UI Electron/Vite \\
\addlinespace
loudness & Orchestrateur & Contrôle du volume système (monter, baisser, régler à X\%) \\
\addlinespace
brightness & Orchestrateur & Contrôle de la luminosité de l'écran \\
\addlinespace
node-window-manager & Orchestrateur & Liste/focus/minimize/maximize les fenêtres Windows (optionnel, natif) \\
\addlinespace
active-win & Orchestrateur & Récupère la fenêtre active courante \\
\addlinespace
@nut-tree-fork/nut-js & Orchestrateur & Simulation clavier/souris (taper du texte, bouger la souris, cliquer) \\
\addlinespace
child\_process (Node intégré) & Orchestrateur & Lance/ferme des applications système (\texttt{start}, \texttt{taskkill}, \texttt{calc}, \texttt{chrome}…) \\
\addlinespace
node:fs (Node intégré) & Orchestrateur & Opérations fichiers (créer dossier, lire, écrire, déplacer, supprimer) \\
\addlinespace
FastAPI & IA (Python) & Serveur HTTP Python exposant les endpoints IA (\texttt{/health}, \texttt{/ask}…) \\
\addlinespace
OpenAI GPT-4o & IA (Python) & LLM : repli \emph{optionnel} (activé via \texttt{USE\_LLM=true}) pour les intentions très complexes \\
\addlinespace
spaCy & IA (Python) & NLP local : NER, POS tagging, dependency parsing, priorité, participants \\
\addlinespace
TextBlob & IA (Python) & Analyse de sentiment locale (polarity + subjectivity) \\
\addlinespace
ChromaDB & IA (Python) & Base de données vectorielle = mémoire long-terme de l'assistant \\
\addlinespace
Whisper (STT) & IA (Python) & Transcription voix → texte (Speech-to-Text, offline) \\
\addlinespace
TTS & IA (Python) & Synthèse vocale texte → audio (réponse parlée) \\
\addlinespace
Electron & UI & Encapsule l'app React dans une fenêtre desktop native \\
\addlinespace
React + Vite & UI & Interface utilisateur (Dashboard, Avatar, Historique, Réglages) \\
\addlinespace
TailwindCSS & UI & Styles CSS utilitaires \\
\addlinespace
Framer Motion & UI & Animations (avatar LED, transitions) \\
\addlinespace
concurrently & Monorepo & Lance les 3 couches en parallèle avec \texttt{npm run dev} \\

\bottomrule
\end{longtable}

% ─────────────────────────────────────────────
\section{Modifications apportées}

\begin{longtable}{>{\ttfamily}p{7cm} p{8cm}}
\toprule
\normalfont\bfseries Fichier & \bfseries Ce qui a changé \\
\midrule
\endhead

apps/orchestrator/src/desktop/apps.js
  & Table d'alias français/anglais (\texttt{chrome}, \texttt{calculatrice}, \texttt{google}…) → résolution vers le bon exécutable + nettoyage des articles (« la », « le »…) + repli si \texttt{start} échoue. \\
\addlinespace
apps/orchestrator/src/router/commandRouter.js
  & Regex plus larges : accepte l'anglais (\texttt{open}, \texttt{start}, \texttt{launch}) + fautes de frappe (\texttt{ouvrire}, \texttt{augmanter}, \texttt{lumenisite}…) + \texttt{volume.set} / \texttt{brightness.set} déclenchés dès qu'un nombre est présent. \\
\addlinespace
apps/orchestrator/src/index.js
  & Gestion claire du port occupé (\texttt{EADDRINUSE}) + handlers globaux pour éviter le crash/déconnexion automatique sur erreur ponctuelle. \\

\bottomrule
\end{longtable}

% ─────────────────────────────────────────────
\section{Flux d'une commande}

\subsection*{Commande locale (offline)}

\begin{lstlisting}
Utilisateur tape : "ouvre chrome"
  -> Orchestrateur : regex match -> app.open { name: "chrome" }
  -> resolveApp("chrome") -> "chrome" (alias Win)
  -> safeExec('start "" "chrome"') -> Chrome s'ouvre
  -> Reponse : "Ouverture de chrome."
\end{lstlisting}

\subsection*{Commande complexe (NLP local)}

\begin{lstlisting}
Utilisateur tape : "appelle Alice demain c'est urgent"
  -> Orchestrateur -> POST http://localhost:8000/api/analyze
  -> Pipeline NLP : NLU + spaCy + TextBlob
  -> intention + entites (Alice) + priorite (5) + echeance (demain)
  -> reponse generee localement (sans LLM)
  -> UI : badge d'intention + reponse
\end{lstlisting}

% ─────────────────────────────────────────────
\section{En résumé}

\textbf{La compréhension du langage est désormais 100\% locale (NLP).}
Le pipeline spaCy + TextBlob analyse l'intention, les entités, la priorité
et le sentiment sans appel réseau ni clé API. Le contrôle du bureau (volume,
luminosité, applications, fichiers) reste \textbf{100\% offline} via
l'orchestrateur Node.js. GPT-4o n'est plus qu'un \textbf{repli optionnel}.

\end{document}
